\subsection{Reproducibility}

\paraib{Accessing the artifact(s)}
%
Our artifacts are open-source, and publicly available on GitHub (Puffer, Mahimahi, and CellReplay). The ABR algorithms that we used are also included in Puffer's repository. They can be found in the "abr" directory, and models such as Pensieve, Gelato, and TTP can be accessed in the "third\_party" directory. Users can clone each project repository, apply the patches described in \ref{s:artifact_abstract}, and configure each artifact following their documentation.

\paraib{Experiment Workflow}
%
To conduct an experiment, the user has to initialize Puffer's web client and media server, and in another terminal start a Mahimahi or CellReplay shell using the desired trace files. Inside the emulated namespace, launch a headless browser instance and connect to the client to start a streaming session using an automated script. In a separate terminal, use tcpdump to capture packets on port 50001, which is the default port where Puffer sends video segments from the media server to the web client. After the test run, export Puffer's logs and the captured TCP packets for analysis. Our methodology is explained in more detail in Section \ref{s:design}.

\paraib{Evaluation and Expected Result}
%
When comparing emulators, we did expect CellReplay to perform better due to its advanced design compared to Mahimahi. We observed some changes in RTT and delivery rate between emulators, but we also expected the SSIM and BufRatio rankings to change slightly. Instead, they remained constant with only a slight improvement when CellReplay is used. SSIM and BufRatio patterns for each ABR algorithm were nearly identical in Mahimahi and CellReplay when the same trace was used. There were slight deviations in the patterns that became clear when we took the mean values for each metric across all test runs.

\paraib{Experiment Customization}
%
For our ABR algorithm configurations, we used \texttt{cubic} congestion control algorithm, which is default in Ubuntu, but if desired, users can also use \texttt{bbr}, by modifying settings.yml file. In addition, we used the recommended configuration for each ABR algorithm, but for custom experiments, these configurations can be modified.


%%% Local Variables:
%%% mode: latex
%%% TeX-master: "../thesis"
%%% End:
